\section{Method}
\label{sec:method}

\begin{figure}[t]
\centering
\begin{tikzpicture}[
  node distance=1.0cm and 0.7cm,
  box/.style={draw, rounded corners=3pt, fill=blue!8, minimum width=2.2cm,
               minimum height=0.65cm, align=center, font=\small},
  dbox/.style={draw, dashed, rounded corners=3pt, fill=gray!10, minimum width=2.2cm,
               minimum height=0.65cm, align=center, font=\small},
  arrow/.style={-{Stealth[length=5pt]}, thick},
  label/.style={font=\scriptsize, align=center}
]

% Row 1: Input frames
\node[box] (frames) {Input frames\\$x_t, x_{t+1}, \ldots$};

% Row 2: Encoder
\node[box, below=of frames] (encoder) {ViT Encoder $f_\theta$\\(frozen LeWM)};
\draw[arrow] (frames) -- (encoder);

% Split: z_t and z_{t+1}
\node[box, below left=0.8cm and 0.5cm of encoder] (zt) {$z_t = f_\theta(x_t)$};
\node[box, below right=0.8cm and 0.5cm of encoder] (zt1) {$z_{t+1} = f_\theta(x_{t+1})$};
\draw[arrow] (encoder) -- (zt);
\draw[arrow] (encoder) -- (zt1);

% Predictor
\node[box, below=0.8cm of zt] (pred) {JEPA Predictor $g_\phi$\\(frozen, zero-action)};
\draw[arrow] (zt) -- (pred);

% Predicted z
\node[box, below=of pred] (zhat) {$\hat{z}_{t+1} = g_\phi(z_t)$};
\draw[arrow] (pred) -- (zhat);

% LPE
\node[box, below right=0.8cm and 0.0cm of zhat, xshift=0.7cm] (lpe)
  {LPE: $d_{e,t} = \|\hat{z}_{t+1} - z_{t+1}\|_2^2$};
\draw[arrow] (zhat) -- (lpe);
\draw[arrow] (zt1) -- (lpe);

% Aggregation
\node[box, below=of lpe] (agg)
  {Episode aggregation\\$S_e = A(\{d_{e,t}\})$};
\draw[arrow] (lpe) -- (agg);

% Calibration node (side)
\node[dbox, right=1.4cm of agg] (cal)
  {Threshold $\tau$\\(p95, normal-only)};

% Decision
\node[box, below=of agg] (dec)
  {Decision: $\hat{y}_e = \mathbf{1}[S_e \geq \tau]$};
\draw[arrow] (agg) -- (dec);
\draw[arrow] (cal) -- (dec);

\end{tikzpicture}
\caption{%
LeWM-Glitch scoring pipeline. Gameplay frames are encoded by a frozen ViT-based
LeWorldModel~\cite{lewm} encoder. A frozen JEPA predictor estimates the next latent state.
The latent prediction error (LPE, Eq.~\ref{eq:lpe}) is aggregated over windows to an episode
score $S_e$ (Eq.~\ref{eq:aggregation}). The detection threshold $\tau$ is set from
normal-only calibration episodes (Eq.~\ref{eq:threshold}); no buggy episodes are used at any
stage before evaluation. This pipeline does not perform temporal localization.%
}
\label{fig:pipeline}
\end{figure}

\subsection{JEPA-Style Latent-Surprise Detector}

LeWM-Glitch is built on the frozen LeWorldModel (LeWM)~\cite{lewm} checkpoint --- a
JEPA-style end-to-end pixel-conditioned dynamics model with a ViT encoder, an AdaLN-conditioned
predictor, and a SIGReg spectral regularizer. Given a window of consecutive gameplay frames, the
encoder $f_\theta$ maps each frame to a latent tensor; the JEPA predictor $g_\phi$ estimates
the next latent state from the preceding context. The \emph{latent prediction error} (LPE) per
window follows Eq.~\ref{eq:lpe}. We evaluate the following score variants:
\begin{equation}
d_{e,t} \in \{\texttt{mse\_mean},\,\texttt{mse\_max},\,\texttt{mse\_top2\_mean},
  \,\texttt{l2\_mean},\,\texttt{l2\_max},\,\texttt{l2\_top2\_mean}\}.
\end{equation}
These differ only in the elementwise reduction (MSE vs.~L2) and the within-window
summary; all are derived from the same frozen checkpoint.

The TempGlitch evaluation path uses a documented zero-action input because verified action
traces are not available in the current dataset adapter. The R5-XGame path uses its own
action-present protocol. Because the two evaluation families differ in domain and action mode,
no cross-family comparison is used to argue for action-conditioning benefit.

\subsection{Episode-Level Aggregation}

Window scores are aggregated to a single episode score $S_e$ as in Eq.~\ref{eq:aggregation}.
\texttt{mean} aggregation emphasizes sustained elevated surprise across the episode;
\texttt{max} and \texttt{top2\_mean} are more sensitive to brief, high-intensity anomalies.
Results are reported per configuration; no post-hoc selection on evaluation data is performed.

\subsection{Baseline Families}

Two simple baselines require no learned components: \texttt{frame\_diff} measures
frame-to-frame visual change; \texttt{feature\_distance} fits a normal reference on
train-normal data and scores Euclidean distance from that reference at test time.
The exact-support K1 comparison adds three learned baselines trained only on allowed
train-normal data: a convolutional LSTM (\texttt{cnn\_lstm}), a video autoencoder
(\texttt{video\_autoencoder}), and a video transformer (\texttt{video\_transformer}).
All baselines use the same episode-level aggregation operators as LeWM.

No fully fine-tuned TimeSformer or VideoMAE anomaly detector was included because the available
train-normal support (fewer than 50 episodes per split) is insufficient to train such models
without severe overfitting. We acknowledge this as a limitation.

\subsection{Normal-Only Threshold Calibration}

For every configuration, $\tau$ is the empirical 95th-percentile episode score over designated
calibration-normal episodes (Eq.~\ref{eq:threshold}). This is an operating-point rule;
buggy episodes are not used at any stage before evaluation. Locked test data is neither
materialized nor scored~\cite{kaufman2011leakage}.
